\documentclass[11pt]{article}
\usepackage[margin=1in]{geometry}
\usepackage{amsmath,amssymb,amsthm}
\usepackage{enumitem}
\usepackage{tikz}
\usetikzlibrary{positioning,arrows.meta,calc}
\usepackage[backend=biber,style=apa,natbib]{biblatex}
\addbibresource{problem.bib}
\usepackage{hyperref}
\usepackage{xcolor}
\hypersetup{colorlinks=true,linkcolor=blue!50!black,urlcolor=blue!50!black,citecolor=blue!50!black}

\newtheorem{definition}{Definition}
\newtheorem{proposition}{Proposition}
\newtheorem{theorem}{Theorem}
\newtheorem{corollary}{Corollary}
\newtheorem{remark}{Remark}
\newtheorem{example}{Example}

\title{elaboration adjunctions: constructive retraction and base-shift naturality}
\author{T.\ Savich, Indiana University\\
\small Companion paper to LX.tex; closes Open Problems 1 and 4}
\date{\today}

\begin{document}
\maketitle

\noindent\fbox{\parbox{\textwidth}{\small\textbf{Status:} Working
document. The two propositions below are stated and proved at a
level adequate for working out the structure; specific cases
(RMB-LX-to-CO under base shift) need to be checked computationally
before either is published as a result. The retraction construction
of \S\ref{sec:retraction} requires a definability condition on
elaborated primitives; this condition is satisfied for the worked
examples but is itself a hypothesis about other cases.}}
\bigskip

\begin{abstract}
We close two of the open problems posed in LX.tex. First, the
retraction $r$ of an elaboration adjunction admits an explicit construction
in the case where the elaborated practice adds primitives that are
definable in the base practice; the unique-up-to-natural-isomorphism
condition reduces to definability of the added primitives. Second,
under base shift, elaboration adjunctions are preserved: if a primitive
signature is base-parametric and its cost function is structurally
defined, then the base-shift functor $\Phi_{b \to b'}$ commutes with
both the elaboration embedding $\iota$ and the retraction $r$, so the
adjunction transports cleanly across bases. Together these results
turn the LX framework from a proposal into a structure whose
manipulations are checkable.
\end{abstract}

\section{Setting}

We assume the framework of LX.tex \S\S 1--4: a typed primitive
signature $\Sigma$, the free coloured operad $\mathbf{Op}(\Sigma)$ of
programs over it, a metaphor labelling $\mu : \Sigma \to 2^\Pi$
inducing sub-operads $\mathbf{Op}_\pi(\Sigma)$ for each
$\pi \in \Pi$, and the elaboration adjunction $r \dashv \iota$ between
$\mathbf{Op}_{\pi_1}$ (the core) and $\mathbf{Op}_{\pi_2}$ (the
elaborated practice) for a chosen pair $\pi_1, \pi_2$ of metaphor
labels.

For simplicity we restrict to the case where the elaborated metaphor
$\pi_2$ has the property
\[
\Sigma^{\pi_2} \;=\; \Sigma^{\pi_1} \;\cup\; \Delta\Sigma
\]
for some non-empty set $\Delta\Sigma$ of new primitives. We call
this the \emph{primitive-extension case}. RMB extends CO by adding
$\mathrm{delta\_to\_next\_base}_b$; FCS extends PFS by adding the
recursive composition operator. Both are primitive-extension cases.

\section{Constructive retraction}\label{sec:retraction}

The retraction $r : \mathbf{Op}_{\pi_2} \to \mathbf{Op}_{\pi_1}$
required by Definition~1 of LX.tex is intuitively ``strip the new
primitives.'' Stripping is only well-defined if each new primitive
admits a definition in terms of the old ones.

\begin{definition}[Definability]\label{def:definable}
A primitive $\sigma \in \Delta\Sigma$ is \emph{definable in
$\Sigma^{\pi_1}$} if there exists a finite program
$P_\sigma \in \mathbf{Op}(\Sigma^{\pi_1})$ of the same arity and
output type as $\sigma$, such that
$\llbracket P_\sigma \rrbracket = \llbracket \sigma \rrbracket$
as partial functions on the semantic universe.
\end{definition}

\begin{example}
$\mathrm{delta\_to\_next\_base}_b(x)$ is definable in
$\Sigma^{\mathrm{CO}}$ as the program that iterates
$\mathrm{succ}$ from $x$ until reaching a multiple of $b$ and returns
the iteration count. The definition uses only successor and equality
testing; both are in $\Sigma^{\mathrm{CO}}$.
\end{example}

\begin{theorem}[Constructive retraction]\label{thm:retraction}
Suppose $\Sigma^{\pi_2} = \Sigma^{\pi_1} \cup \Delta\Sigma$ and every
$\sigma \in \Delta\Sigma$ is definable in $\Sigma^{\pi_1}$, with
chosen definitions $\{P_\sigma\}_{\sigma \in \Delta\Sigma}$. Define
$r : \mathbf{Op}_{\pi_2} \to \mathbf{Op}_{\pi_1}$ on operations by
operadic substitution: replace every occurrence of $\sigma \in
\Delta\Sigma$ in any tree by its definition $P_\sigma$.

Then $r$ is well-defined, satisfies $r \circ \iota =
\mathrm{id}_{\mathbf{Op}_{\pi_1}}$, and is unique up to natural
isomorphism given the chosen definitions $\{P_\sigma\}$.
\end{theorem}

\begin{proof}[Proof sketch]
\emph{Well-defined.} Each $\sigma \in \Delta\Sigma$ has finite
definition $P_\sigma$; operadic substitution of $P_\sigma$ for
$\sigma$ in any program produces a finite program in
$\mathbf{Op}(\Sigma^{\pi_1})$. Since $\Sigma^{\pi_1}$ is closed under
operadic substitution within its own label (Prop.~1 of LX.tex),
the result lies in $\mathbf{Op}_{\pi_1}$.

\emph{Retracts $\iota$.} The embedding $\iota$ takes a
$\pi_1$-program $\alpha$ to its image in $\mathbf{Op}_{\pi_2}$, which
is structurally the same tree (no $\Delta\Sigma$ primitive
appears). Applying $r$ then finds no $\Delta\Sigma$ primitive to
substitute, so $r(\iota(\alpha)) = \alpha$.

\emph{Unique up to natural iso.} Suppose $r'$ is another retraction
arising from a different choice of definitions
$\{P'_\sigma\}_{\sigma \in \Delta\Sigma}$. Each $P_\sigma$ and
$P'_\sigma$ compute the same partial function, so they are
semantically equivalent. The natural isomorphism $\nu : r \Rightarrow
r'$ sends each $r(\alpha)$ to $r'(\alpha)$ via the semantic
equivalence; naturality follows from compositionality of substitution.
\end{proof}

\begin{corollary}[Adjointness reduces to definability]
In the primitive-extension case, the elaboration adjunction $r \dashv \iota$
exists if and only if every $\sigma \in \Delta\Sigma$ is definable in
$\Sigma^{\pi_1}$.
\end{corollary}

\begin{remark}[Why definability is exactly the right condition]
The Brandomian intuition behind LX is that the elaborated vocabulary
$V_{\pi_2}$ makes \emph{explicit} something already implicit in the
practice of using $V_{\pi_1}$. Definability is the categorical
reading of ``already implicit'': the new primitive $\sigma$ is
implicit in $\Sigma^{\pi_1}$ when it can be defined from existing
primitives — even if defining it explicitly was not part of the
$\pi_1$-practice's expressive resources. The retraction $r$ makes
the definition manifest by unfolding it.
\end{remark}

\begin{remark}[The non-primitive-extension case]
Some LX relations are not of primitive-extension form. The
elaborated practice might revise the semantics of an existing
primitive, or replace one primitive with two specialised variants.
For these cases Theorem~\ref{thm:retraction} does not apply
directly; a more delicate analysis is needed. The simplest
counter-construction is one where the elaborated practice changes
the input type signature, in which case the embedding $\iota$ is no
longer a strict inclusion. We leave this for future work.
\end{remark}

\section{Multi-base preservation of elaboration adjunctions}\label{sec:multi-base}

We now show that under base shift, elaboration adjunctions transport along
the obvious renaming functor. This requires the cost function to be
structurally defined in the sense of ALGEBRAIC.tex \S2.

Recall the base-shift functor $\Phi_{b \to b'} :
\mathbf{Op}(\Sigma_b) \to \mathbf{Op}(\Sigma_{b'})$ which renames
base-tagged primitives (e.g.\ $\mathrm{delta\_to\_next\_base}_b
\mapsto \mathrm{delta\_to\_next\_base}_{b'}$) and is the identity on
all base-independent primitives.

\begin{theorem}[Multi-base LX preservation]\label{thm:multi-base-lx}
Let $r_b \dashv \iota_b$ be an elaboration adjunction between
$\mathbf{Op}_{\pi_1, b}$ and $\mathbf{Op}_{\pi_2, b}$, where both
$\pi_1$ and $\pi_2$ are primitive-extension cases over a
base-parametric signature $\Sigma$. If the cost function is
structurally defined, then there is a corresponding elaboration adjunction
$r_{b'} \dashv \iota_{b'}$ in base $b'$, and the diagrams
\begin{align}
\Phi_{b \to b'} \circ \iota_b &= \iota_{b'} \circ \Phi_{b \to b'}, \\
\Phi_{b \to b'} \circ r_b &= r_{b'} \circ \Phi_{b \to b'}
\end{align}
commute.
\end{theorem}

\begin{proof}[Proof sketch]
\emph{Construction of $\iota_{b'}$ and $r_{b'}$.} Both sub-operads
$\mathbf{Op}_{\pi_1, b'}$ and $\mathbf{Op}_{\pi_2, b'}$ exist by
construction (the metaphor labelling is base-independent and the
primitive signature is base-parametric). The embedding $\iota_{b'}$
is again the structural identity-embedding. The retraction $r_{b'}$
is constructed by Theorem~\ref{thm:retraction} provided each
$\sigma \in \Delta\Sigma_{b'}$ is definable in $\Sigma^{\pi_1}_{b'}$.

\emph{Definability transports.} For each base-tagged primitive
$\sigma_b \in \Delta\Sigma_b$, the definition $P_{\sigma_b}$ in
$\Sigma^{\pi_1}_b$ uses only base-independent primitives plus
possibly base-tagged primitives in $\Sigma^{\pi_1}_b$. The
re-encoding $\Phi_{b \to b'}(P_{\sigma_b}) =: P_{\sigma_{b'}}$ is
a valid definition of $\sigma_{b'}$ in $\Sigma^{\pi_1}_{b'}$
because the structural cost condition guarantees the semantics of
all base-tagged primitives is the same up to substituting $b$ for
$b'$ in mod-arithmetic. So definability transports.

\emph{$\iota$ commutes with $\Phi$.} The embedding $\iota_b$ is the
identity on tree structure; $\Phi_{b \to b'}$ renames only
base-tagged labels at the primitive level. So
$\Phi(\iota_b(\alpha)) = \iota_{b'}(\Phi(\alpha))$ tree-by-tree.

\emph{$r$ commutes with $\Phi$.} The retraction $r_b$ substitutes
each $\sigma_b \in \Delta\Sigma_b$ by $P_{\sigma_b}$. The retraction
$r_{b'}$ substitutes each $\sigma_{b'} \in \Delta\Sigma_{b'}$ by
$\Phi(P_{\sigma_b})$. Since $\Phi$ commutes with substitution
(operadic substitution and renaming are independent operations),
$\Phi(r_b(\alpha)) = r_{b'}(\Phi(\alpha))$.

\emph{Adjointness preserves.} The universal property of $r_b
\dashv \iota_b$ is a statement about morphisms in
$\mathbf{Op}_{\pi_2, b}$; since $\Phi$ is an isomorphism of
operads, it preserves the universal property. Therefore
$r_{b'} \dashv \iota_{b'}$.
\end{proof}

\begin{corollary}[The unit is base-natural]
The unit $\eta^b : \mathrm{id} \Rightarrow \iota_b \circ r_b$ of the
elaboration adjunction in base $b$ corresponds under $\Phi_{b \to b'}$ to the
unit $\eta^{b'}$ of the elaboration adjunction in base $b'$. Formally,
$\Phi_{b \to b'}(\eta^b_\alpha) = \eta^{b'}_{\Phi(\alpha)}$ for all
$\alpha \in \mathbf{Op}_{\pi_2, b}$.
\end{corollary}

\section{Worked example: RMB LX to CO across bases}

Consider the elaboration of CO into RMB in base $b$. Let:
\begin{itemize}[noitemsep,topsep=2pt]
  \item $\Sigma^{\mathrm{CO}}_b = \{\mathrm{succ},\
        \mathrm{pred},\ \mathrm{iter\_succ},\ \mathrm{iter\_pred}\}$
  \item $\Sigma^{\mathrm{RMB}}_b = \Sigma^{\mathrm{CO}}_b \cup
        \Delta\Sigma_b$, where
        $\Delta\Sigma_b = \{\mathrm{delta\_to\_next\_base}_b,\
        \mathrm{delta\_from\_prev\_base}_b\}$.
\end{itemize}
Each base-tagged primitive is definable in $\Sigma^{\mathrm{CO}}_b$
by counting:
\begin{align}
P_{\mathrm{delta\_to\_next\_base}_b}(x) &= \ \text{iterate } \mathrm{succ}
  \text{ from } x \text{ until next multiple of } b\text{; return count;} \\
P_{\mathrm{delta\_from\_prev\_base}_b}(x) &= \ \text{iterate }
  \mathrm{pred} \text{ from } x \text{ until previous multiple of }
  b\text{; return count.}
\end{align}
Both definitions are base-parametric: their structure is the same
in any base $b$, with only the modulus changing.

By Theorem~\ref{thm:retraction}, $r_b : \mathbf{Op}_{\mathrm{RMB}, b}
\to \mathbf{Op}_{\mathrm{CO}, b}$ exists and is unique up to natural
isomorphism. By Theorem~\ref{thm:multi-base-lx}, this adjunction
transports across bases: the RMB-LX-CO adjunction in base $10$
corresponds, under $\Phi_{10 \to 12}$, to the RMB-LX-CO adjunction
in base $12$.

\paragraph{Empirical check.} The synthesis prototype at
\texttt{formalization/synthesis/synth\_lazy.pl} now accepts a base
parameter. Running the make-base-subtraction target under bases $5$,
$7$, $10$, and $12$ on structurally-matched inputs (each subtrahend
chosen to share the same $\delta_b$-profile across bases) should
produce the same canonical four-step program shape: $\delta$ +
$\mathrm{iter\_succ}$ + $\mathrm{iter\_succ}$ +
$\mathrm{iter\_pred\_aligned}$. Theorem~\ref{thm:multi-base-lx}
predicts this; the structurally-defined cost function in the
prototype is the condition that makes the prediction follow without
empirical surprise.

\section{Open questions remaining after this paper}

\begin{enumerate}[noitemsep,topsep=2pt,label=(\arabic*)]
  \item \textbf{Non-primitive-extension LX.} The constructive
        retraction requires the elaborated practice to add new
        primitives over the core. When the elaboration revises
        existing primitives or changes the type signature, a more
        general construction is needed. What is the right
        generalisation?
  \item \textbf{Multiple definitions, multiple retractions.}
        Theorem~\ref{thm:retraction} fixes a choice of
        $\{P_\sigma\}$ and gives uniqueness up to natural
        isomorphism. Are there cases where different choices of
        definition produce genuinely different elaboration adjunctions —
        differing in unit, in cost, in elaboration depth?
  \item \textbf{Descent along the metaphor cover.} Theorem~2 of
        LX.tex (LX-towers compose) holds for primitive-extension
        cases by Theorem~\ref{thm:retraction}. But the empirical
        elaboration chains in the project pass through metaphor
        labels in ways that are not always primitive-extensions.
        Whether the full chain CO $\to$ RMB $\to$ Chunking $\to$ Skip
        Counting $\to \ldots \to$ CGOB Division is an LX-tower of
        primitive-extension steps, or whether some step is
        non-extension, is an empirical question to settle.
  \item \textbf{Computability of the unit.} Open Problem 5 of
        LX.tex asked whether the explicitness witness $\eta_\alpha$
        is algebraically computable. Under
        Theorem~\ref{thm:retraction}, the answer is yes for
        primitive-extension cases: $\eta_\alpha$ is the morphism
        $\alpha \to \iota(r(\alpha))$ in $\mathbf{Op}_{\pi_2}$, and
        $\iota(r(\alpha))$ is computed by substituting each
        $\sigma \in \Delta\Sigma$ in $\alpha$ by $P_\sigma$ and
        then re-embedding. This is a definable derivation, so
        $\eta_\alpha$ can be read off mechanically from $\alpha$
        and the chosen definitions $\{P_\sigma\}$. This addresses
        OP5 in the affirmative for the primitive-extension setting.
\end{enumerate}

\section{Note on what this paper does and does not do}

These two results turn LX-as-adjunction from a proposal into a
checkable structure in the primitive-extension case. They do
\emph{not} extend the proposal to all of Brandom's documented LX
relations: some elaborations may not be primitive-extension form,
and for those the constructive retraction may not exist. They also
do not validate the claim that the empirical strategy chain in the
project's arithmetic registry is a chain of elaboration adjunctions in this
specific sense; that remains an empirical question per strategy.

What is now verifiable is the per-case claim that a specific
elaboration (RMB-from-CO, FCS-from-PFS, possibly Chunking-from-RMB)
satisfies the primitive-extension condition with definable new
primitives, so that the elaboration adjunction structure follows from
Theorem~\ref{thm:retraction}. The prototype's multi-base run
provides the empirical check on Theorem~\ref{thm:multi-base-lx} for
RMB-LX-CO specifically.

\printbibliography

\end{document}
